\documentclass[a4paper,12pt]{article}
\input{ReportHeader}
\graphicspath{{../Images/}}

\begin{document}
%----------------------------------------------------------------------------------------
%	TITLE+
%----------------------------------------------------------------------------------------
% Title Page
\begin{titlepage}
    \centering
    \vspace*{2cm}
    \Huge{\textbf{EEE3027 Integrated Circuit Design}}\\[0.5cm]
    \Large{\textbf{Semester 1 Report}}\\ 
    \Large{Sahas Talasila \textit{230057896}}
    \vfill
\end{titlepage}

% Table of Contents
\tableofcontents
\newpage

\begin{abstract}
    This is a placeholder for the abstract of the report.
\end{abstract}

\section{Introduction}

\section{NMOS and PMOS I-V Characteristics Lab 1}

\subsection{NMOS Characteristics}

Include theory on how a transistor works, regions of operation, and also beta.

Explain the differences between holes and electrons, and explain why it will be useful to know both characteristics, 
and how it shows tthis in an NMOS

Explain the code used to generate the plots,

Explain the plots.

Finally, show the calculations for each curve.

\subsection{PMOS characteristics}

Have already included theory on how a transistor works, regions of operation, and also beta.

Explain the code to generate the plots (just a small difference),

Explain the plots. (plot of Id vs Vds) 

Finally, show the calculations for each curve. Example shown below

\begin{itemize}
  \item \textbf{EXPECTED CURVES:}
  \begin{itemize}
    \item \textbf{Curve 1 ($V_{GS} = 0.4\,\text{V}$):}
      \begin{itemize}
        \item Barely turns on ($V_{GS} = V_{th}$)
        \item Very low current
        \item Saturation at $I_{DS} \approx 0\,\text{mA}$
      \end{itemize}

    \item \textbf{Curve 2 ($V_{GS} = 0.6\,\text{V}$):}
      \begin{itemize}
        \item $V_{DSsat} = 0.6 - 0.4 = 0.2\,\text{V}$
        \item Saturation current (approx) $= 0.04\,\text{mA}$
        \item $I_{DS,sat} = \frac{\beta}{2}(V_{GS} - V_{th})^2 = 1\,\text{m} \times (0.2)^2 = 0.04\,\text{mA}$
      \end{itemize}

    \item \textbf{Curve 3 ($V_{GS} = 0.8\,\text{V}$):}
      \begin{itemize}
        \item $V_{DSsat} = 0.8 - 0.4 = 0.4\,\text{V}$
        \item Saturation current (approx) $= 0.16\,\text{mA}$
        \item $I_{DS,sat} = 1\,\text{m} \times (0.4)^2 = 0.16\,\text{mA}$
      \end{itemize}

    \item \textbf{Curve 4 ($V_{GS} = 1.0\,\text{V}$):}
      \begin{itemize}
        \item $V_{DSsat} = 1.0 - 0.4 = 0.6\,\text{V}$
        \item Saturation current (approx) $= 0.36\,\text{mA}$
        \item $I_{DS,sat} = 1\,\text{m} \times (0.6)^2 = 0.36\,\text{mA}$
      \end{itemize}

    \item \textbf{Curve 5 ($V_{GS} = 1.2\,\text{V}$):}
      \begin{itemize}
        \item $V_{DSsat} = 1.2 - 0.4 = 0.8\,\text{V}$
        \item Saturation current (approx) $= 0.64\,\text{mA}$
        \item $I_{DS,sat} = 1\,\text{m} \times (0.8)^2 = 0.64\,\text{mA}$
      \end{itemize}

    \item \textbf{Each curve:}
      \begin{itemize}
        \item Linear region: $0$ to $(V_{GS} - V_{th})$
        \item Saturation region: $(V_{GS} - V_{th})$ to $1.2\,\text{V}$
        \item Higher $V_{GS} \Rightarrow$ higher saturation current
        \item Saturation current $\propto (V_{GS} - V_{th})^2$
      \end{itemize}
  \end{itemize}
\end{itemize}


\subsection{Beta Calculation}

Explain what beta is, and how to calculate it.
Show the calculations for both NMOS and PMOS.

Explain the code again and then show a plot of beta and then gradient calculation.

Explain what the gradient means.

Show image with gradients, then compare theoretical and experimental values.

Explain the relationship between mu, Cox, W and L, and how changing W and L will affect beta.

Show plot with l times 2 and how gradient has halved, refering back to the beta equation and the equivalent versions.

Show how VGS also affects beta. (has to be in saturation region for this to be true), otherwise it is not valid.
Show plot with VGS changed and how beta changes accordingly.
Refer back to beta equation again.

Beta vs width or length table

\subsection{How SPICE calculates non linear equations}
Include brief theory on non linear equations in circuits.

Add diagram of diode and resistor in series.

Explain how SPICE uses piecewise linear approximation to solve non linear equations.

And explain how newton raphson method is used to converge to a solution for the non linear equations

Show worked example of diode equation and how it is approximated using piecewise linear approximation.

Refer to python code and also the cobweb diagram. 

Compare bisection method to newton raphson method and explain why newton raphson is better for SPICE.

Add flowchart//pseudocode of how SPICE solves non linear equations using newton raphson method.

\section{Lab 1 part 2: CMOS Inverter and Experiments}

\subsection{CMOS Inverter Diagram and Theory}

Explain the CMOS inverter circuit diagram and how a CMOS inverter works.

Explain the components used (NMOS and PMOS) and how they work together to create the inverter.

\subsection{CMOS Missing SPICE}

Explain the missing SPICE code and how it works.

Explain the additions and annotations made to the original code.

\subsection{Transfer Curves}

Explain what the transfer curves and how they were generated.

Point out the plots and explain the different regions of operation for PMOS and NMOS.

Explain when NMOS is on and when PMOS is on/triode/saturation.

Analysis table showing what happens at each voltage level. 

\subsection{NMOS Pull Up}

Explain the modifications made to the CMOS inverter to create an NMOS pull up.

Show the SPICE code used.

Explain the transfer curves generated and compare to the CMOS inverter.

Show the differences and why they occur.

Show the calculations for the threshold voltages of both NMOS and PMOS in the circuit, 
and explain how they were calculated and how they affect the operation of the inverter.

Show how the NMOS pull up is worse than the CMOS inverter.

Answer Questions from lab brief.

\subsection{CMOS Inverter Design}

Explain the requirements for the CMOS inverter design and what they mean. 

\end{document}